\documentclass[11pt,a4paper]{article}
\usepackage{amsmath,amssymb,amsthm}
\usepackage{geometry}
\geometry{margin=1in}
\usepackage{hyperref}
\usepackage{enumitem}

\newtheorem{theorem}{Theorem}[section]
\newtheorem{lemma}[theorem]{Lemma}
\newtheorem{proposition}[theorem]{Proposition}
\newtheorem{corollary}[theorem]{Corollary}
\theoremstyle{definition}
\newtheorem{definition}[theorem]{Definition}
\newtheorem{problem}[theorem]{Open Problem}
\theoremstyle{remark}
\newtheorem{remark}[theorem]{Remark}

\title{\bfseries Greedy Dirichlet Sub‑Sum Transform\\
\large A Research Programme on the Selected and Rejected Series}
\author{}
\date{}

\begin{document}
\maketitle

\begin{abstract}
The Greedy Dirichlet Sub‑Sum Transform (GDST) associates to each angle \(\theta\in[0,\pi]\) two complementary Dirichlet series:
\[
E(\theta,s)=\sum_{\delta_n(\theta)=1}\frac{1}{(n+2)^s},\qquad 
\Omega(\theta,s)=\sum_{\delta_n(\theta)=0}\frac{1}{(n+2)^s},
\]
called the \emph{selected} and \emph{rejected} series.
Their sum is \(\zeta(s)-1\).
We outline a research programme that exploits this decomposition to study the Riemann zeta function.
The programme builds on the recently established rigorous foundations of the GDST, replaces the earlier erroneous pole‑based approach with a cancellation and balance mechanism, and formulates precise, accessible open problems.
The ultimate goal is to use the interplay between the selected and rejected series to obtain new insights into the distribution of the zeros of \(\zeta(s)\), potentially culminating in a characterisation of the Riemann Hypothesis as a property of the function \(\theta\mapsto\operatorname{Re}E(\theta,s)\).
\end{abstract}

\section{Introduction}
Let \(\delta_n(\theta)\in\{0,1\}\) be the digits of the greedy harmonic expansion of \(\theta/\pi\) with respect to the harmonic fractions \(\pi/(n+2)\).
The foundational paper~\cite{GDSTfound} proved that for every \(\theta\in[0,\pi]\) the Dirichlet series
\[
E(\theta,s)=\sum_{n=0}^{\infty}\frac{\delta_n(\theta)}{(n+2)^s},\qquad 
\Omega(\theta,s)=\sum_{n=0}^{\infty}\frac{1-\delta_n(\theta)}{(n+2)^s}
\]
converge conditionally and are holomorphic on the half‑plane \(\operatorname{Re}s>0\).
Their sum is the shifted Riemann zeta function:
\begin{equation}\label{eq:sum}
E(\theta,s)+\Omega(\theta,s)=\sum_{n=0}^{\infty}\frac{1}{(n+2)^s}=\zeta(s)-1, \qquad \operatorname{Re}s>0.
\end{equation}
The selected series \(E\) is sparse (its coefficients have natural density zero), while the rejected series \(\Omega\) is co‑sparse (density one) and inherits the pole of \(\zeta(s)\) at \(s=1\).

Two key analytic facts were proved in~\cite{GDSTfound}.
\begin{enumerate}[label=(\arabic*)]
\item On the critical line \(s=\frac12+it\) both series exhibit at most linear growth in \(|t|\).
\item By the identity \eqref{eq:sum} and the convexity bound for \(\zeta(s)\), the sum \(E+\Omega\) is an order of magnitude smaller than the typical size of each term:
\[
E(\theta,\tfrac12+it)=-\Omega(\theta,\tfrac12+it)+O(|t|^{1/6+\varepsilon}) \qquad (|t|\to\infty).
\]
\end{enumerate}
Thus the two series are almost perfect negatives on the critical line.

At a non‑trivial zero \(\rho\) of \(\zeta(s)\) the identity \eqref{eq:sum} becomes
\[
E(\theta,\rho)+\Omega(\theta,\rho)=-1.
\]
Taking real parts yields
\[
\operatorname{Re}E(\theta,\rho)+\operatorname{Re}\Omega(\theta,\rho)=-1.
\]
The Riemann Hypothesis can then be expressed as a property of the selected series alone.

\begin{theorem}[RH equivalence~\cite{GDSTfound}]\label{thm:RH}
All non‑trivial zeros of \(\zeta(s)\) lie on the critical line \(\operatorname{Re}s=\frac12\) if and only if for every \(s\) with \(0<\operatorname{Re}s<1\), \(\operatorname{Re}s\neq\frac12\), and every \(\theta\in[0,\pi]\),
\[
\operatorname{Re}E(\theta,s)\neq -\frac12.
\]
\end{theorem}

The present programme aims to explore the consequences of this reformulation, to deepen the understanding of the selected and rejected series, and to develop tools that might eventually decide the condition in Theorem~\ref{thm:RH}.

\section{Goals}
The programme is divided into four phases.

\begin{enumerate}[label=\textbf{Phase \arabic*},leftmargin=*]
\item \textbf{Analytic structure of \(E\) and \(\Omega\).}
Determine the maximal domain of analyticity for the two series, their growth on vertical lines, and the nature of the balance condition \(\operatorname{Re}E(\theta,s)=-\frac12\) as a function of \(\theta\) and \(s\).
\item \textbf{Correlation and combinatorial structure.}
Express the autocorrelation of the digits \(\delta_n(\theta)\) through the greedy harmonic tree and link it to the value distribution of \(E(\theta,s)\) on the critical line.
\item \textbf{Zeta zeros as balance points.}
Study the set of angles \(\Theta(s)=\{\theta:\operatorname{Re}E(\theta,s)=-\frac12\}\).
Prove existence, uniqueness, and monotonicity properties for \(s\) on the critical line.
Relate the density of \(\Theta(\rho)\) for zeros \(\rho\) to the distribution of the zeros.
\item \textbf{Towards a proof of the RH equivalence.}
Attempt to prove the conjecture of Theorem~\ref{thm:RH} using the combinatorial and analytic machinery developed in Phases~1--3.
\end{enumerate}

\section{Phase 1: Analytic structure}
\subsection{Maximal domain}
From~\cite{GDSTfound} we know that both series converge for \(\operatorname{Re}s>0\).  For a generic \(\theta\) the generating function
\[
G(\theta,z)=\sum_{n=0}^{\infty}\delta_n(\theta)z^{n+2}
\]
has the unit circle as a natural boundary (Fatou--Carlson--Polya).
The corresponding statement for the Dirichlet series is open.

\begin{problem}[Natural boundary for \(E\) and \(\Omega\)]\label{prob:natural}
Prove that for almost all \(\theta\in[0,\pi]\) the line \(\operatorname{Re}s=0\) is a natural boundary for both \(E(\theta,s)\) and \(\Omega(\theta,s)\).
Equivalently, show that the generating function \(G(\theta,z)\) has no analytic continuation beyond the unit disc.
Determine whether there exist exceptional \(\theta\) for which the series continue meromorphically to a larger domain.
\end{problem}

\subsection{Growth on vertical lines}
The bound \(E(\theta,\frac12+it)=O(|t|)\) is coarse.  Tighter estimates would strengthen the cancellation analysis.

\begin{problem}[Improved vertical bounds]\label{prob:growth}
For fixed \(\theta\) investigate the true order of \(E(\theta,\frac12+it)\).
\begin{enumerate}[label=(\alph*)]
\item Prove that \(\int_0^T|E(\theta,\frac12+it)|^2dt \sim C_\theta\,T\log T\) and determine \(C_\theta\).
\item Determine whether \(E(\theta,\frac12+it)=O(|t|^{1/2+\varepsilon})\) or even \(O(|t|^{\varepsilon})\).
\item Relate the exponent to the autocorrelation of \(\delta_n(\theta)\).
\end{enumerate}
\end{problem}

\subsection{The balance function}
For fixed \(s\) with \(\operatorname{Re}s>0\) define the balance function
\[
B_s(\theta)=\operatorname{Re}E(\theta,s)+\frac12,\qquad \theta\in[0,\pi].
\]
The zeros of \(B_s\) are the balance angles.  From the tree representation (\cite{GDSTfound}, Eq.~(5)), \(E(\cdot,s)\) is a regulated function with jumps at the split points of the greedy tree.  Hence \(B_s\) is also regulated.

\begin{problem}[Structure of \(B_s\)]\label{prob:balance}
\begin{enumerate}[label=(\alph*)]
\item For \(s\) in the critical strip, prove that \(B_s\) is of bounded variation if and only if \(\operatorname{Re}s>2\).
\item For \(s=\frac12+it\), study the number of sign changes of \(B_s\) as a function of \(t\).
\item Determine whether \(B_s\) can be expressed as a convolution of the jump measure with a kernel related to the Riemann zeta function.
\end{enumerate}
\end{problem}

\section{Phase 2: Correlation and combinatorial structure}
\subsection{Autocorrelation of the digits}
Define the autocorrelation
\[
R_m(\theta)=\lim_{N\to\infty}\frac{1}{N}\sum_{n=0}^{N-1}\delta_n(\theta)\delta_{n+m}(\theta).
\]
Because the greedy expansion is deterministic but highly irregular, \(R_m\) is expected to decay slowly.
Geere~\cite{Geere2026a} computed approximate correlations; a rigorous theory is lacking.

\begin{problem}[Autocorrelation and the tree]\label{prob:corr}
\begin{enumerate}[label=(\alph*)]
\item Prove the existence of \(R_m(\theta)\) for all \(m\) and almost all \(\theta\).
\item Express \(R_m(\theta)\) using the transition probabilities of the greedy harmonic tree.
\item Link the generating function \(\sum_m R_m(\theta) e^{-2\pi i m t}\) to the spectral measure of the tree's transfer operator.
\end{enumerate}
\end{problem}

\subsection{Value distribution on the critical line}
The mean‑square of \(E(\theta,\frac12+it)\) is controlled by the correlation.

\begin{problem}[Value distribution]\label{prob:dist}
Using the autocorrelation, prove a central limit theorem for \(E(\theta,\frac12+it)\) as \(T\to\infty\) when averaged over \(t\in[0,T]\).
If such a theorem holds, the balance condition \(\operatorname{Re}E=-\frac12\) can be interpreted as a rare event whose probability might be linked to the density of zeros.
\end{problem}

\section{Phase 3: Zeta zeros as balance points}
\subsection{Existence and uniqueness of balance angles}
For a zero \(\rho=\frac12+i\gamma\) we have \(\Theta(\rho)\neq\emptyset\) because \(B_\rho(0)=0+\frac12=\frac12>0\) and \(B_\rho(\pi)=\operatorname{Re}E(\pi,\rho)+\frac12\).  Since \(E(\pi,\rho)=\zeta(\rho)-1-\Omega(\pi,\rho) = -1-\Omega(\pi,\rho)\), we get \(\operatorname{Re}E(\pi,\rho) = -1-\operatorname{Re}\Omega(\pi,\rho) \le -1\) (since \(\operatorname{Re}\Omega\ge 0\)).  Hence \(B_\rho(\pi) \le -\frac12\).  By the intermediate value property (regulated functions have the Darboux property? Regulated functions are not necessarily Darboux, but because \(E\) is a limit of step functions with jumps of shrinking size, it is Darboux).  Thus there exists at least one \(\theta\) with \(B_\rho(\theta)=0\).  Uniqueness is open.

\begin{problem}[Balance angles]\label{prob:angles}
\begin{enumerate}[label=(\alph*)]
\item Prove that for each zero \(\rho\) the equation \(\operatorname{Re}E(\theta,\rho)=-\frac12\) has a unique solution \(\theta_{\mathrm{bal}}(\rho)\).
\item Show that the map \(\rho\mapsto\theta_{\mathrm{bal}}(\rho)\) is strictly increasing with \(\operatorname{Im}\rho\).
\item Determine the asymptotic distribution of \(\{\theta_{\mathrm{bal}}(\rho): 0<\operatorname{Im}\rho\le T\}\) as \(T\to\infty\) and relate it to the \(n\)-point correlation of the zeros.
\end{enumerate}
\end{problem}

\subsection{The balance condition and the explicit formula}
The explicit formula (Weil) expresses sums over zeta zeros in terms of sums over primes.
If one can write \(\operatorname{Re}E(\theta,\frac12+it)\) as a sum over primes with explicit coefficients, the balance condition becomes a prime‑theoretic statement.

\begin{problem}[Explicit formula for the GDST]\label{prob:explicit}
Express \(E(\theta,s)\) as a convolution of the von Mangoldt function with a kernel derived from the greedy tree.
Use this to rewrite the balance condition \(\operatorname{Re}E(\theta,\rho)=-\frac12\) as an identity involving a sum over primes.
\end{problem}

\section{Phase 4: Towards the Riemann Hypothesis}
The ultimate aim is to prove or disprove the equivalence in Theorem~\ref{thm:RH}.
Two distinct strategies emerge from the previous phases.

\subsection{Strategy A: Range of \(B_s\)}
Prove that for \(0<\operatorname{Re}s<1\), \(\operatorname{Re}s\neq\frac12\), the range of \(\theta\mapsto\operatorname{Re}E(\theta,s)\) is an interval (or a finite union of intervals) that does \emph{not} contain \(-\frac12\).
This would require showing that the jumps and the continuous part never allow the value \(-\tfrac12\).

\subsection{Strategy B: Uniqueness of balance and density}
Assuming the existence and uniqueness of \(\theta_{\mathrm{bal}}(\rho)\) (Problem~\ref{prob:angles}), prove that the set \(\{\theta_{\mathrm{bal}}(\rho)\}\) is dense in \([0,\pi]\).
Then show that if a zero off the line existed, its balance angle would violate the density or ordering, leading to a contradiction.

\begin{problem}[RH via balance angles]\label{prob:RHfinal}
Using the results of Problems~\ref{prob:angles}--\ref{prob:explicit}, attempt to prove:
\begin{enumerate}[label=(\alph*)]
\item If \(\rho\) is a zero with \(\operatorname{Re}\rho>\frac12\), then \(\theta_{\mathrm{bal}}(\rho)\) is not realisable (e.g., the function \(B_\rho\) never crosses zero).
\item Alternatively, prove that the balance angle map is a bijection from the set of zeros on the critical line onto a set that is strictly smaller than \([0,\pi]\); an off‑line zero would have to be mapped into the missing set, forcing a contradiction.
\end{enumerate}
\end{problem}

\section{Methodology}
The programme employs a combination of:
\begin{itemize}
\item \textbf{Dirichlet series techniques:} Perron's formula, mean‑value theorems, and the theory of the Mellin transform.
\item \textbf{Combinatorics of the greedy harmonic tree:} split points, cylinder intervals, and transition counts.
\item \textbf{Ergodic theory:} The greedy map is a non‑autonomous dynamical system whose invariant measure controls the statistics of \(\delta_n\).
\item \textbf{Explicit formulas:} Connections to the primes via the von Mangoldt function and the Guinand–Weil explicit formula.
\item \textbf{Numerical exploration:} Large‑scale computation of \(E(\theta,s)\) for moderate heights to test conjectures and guide analytic proofs.
\end{itemize}

\section{Relations to prior work}
The GDST programme corrects and refines the earlier Harmonic Sine Transform (HST) approach of Geere et al.~\cite{HSTprogramme}.
The original claim that the transfer operator \(\mathcal{L}_s\) has determinant \(\zeta(s)/\zeta(2s)\) and that the zeros of \(\zeta\) appear as poles of \(E(\theta,s)\) has been shown to be unsupported.
The present programme abandons the transfer‑operator route and works directly with the selected and rejected Dirichlet series.
The cancellation theorem (Section~1) and the balance formulation (Theorem~\ref{thm:RH}) are the new conceptual pillars.

\section{Open problems summary}
For convenience we collect all open problems stated above, together with the two original problems from~\cite{GDSTfound} concerning the classification of periodic digit sequences and the natural boundary.

\vspace{4pt}
\noindent\textbf{Foundational problems}
\begin{enumerate}[label=\textbf{F\arabic*},leftmargin=*]
\item \textbf{Periodic classification.} Characterise the set of \(\theta\) for which \(\delta_n(\theta)\) is eventually periodic. Prove it is countable and consists of rational multiples of \(\pi\). \hfill (Problem~\ref{prob:periodic} in~\cite{GDSTfound})
\item \textbf{Natural boundary.} Prove that \(\operatorname{Re}s=0\) is a natural boundary for \(E(\theta,s)\) for almost all \(\theta\). \hfill (Problem~\ref{prob:natural})
\end{enumerate}

\noindent\textbf{Phase 1 -- Analytic structure}
\begin{enumerate}[label=\textbf{A\arabic*},leftmargin=*,resume]
\item \textbf{Vertical growth.} Determine the true order of \(E(\theta,\frac12+it)\) and its mean‑square. \hfill (Problem~\ref{prob:growth})
\item \textbf{Balance function.} Study the variation and zeros of \(B_s(\theta)\). \hfill (Problem~\ref{prob:balance})
\end{enumerate}

\noindent\textbf{Phase 2 -- Correlation}
\begin{enumerate}[label=\textbf{C\arabic*},leftmargin=*,resume]
\item \textbf{Autocorrelation.} Prove existence and tree representation of \(R_m(\theta)\). \hfill (Problem~\ref{prob:corr})
\item \textbf{Value distribution.} Establish a central limit theorem for \(E(\theta,\frac12+it)\) in the \(t\)‑aspect. \hfill (Problem~\ref{prob:dist})
\end{enumerate}

\noindent\textbf{Phase 3 -- Zeros as balance points}
\begin{enumerate}[label=\textbf{Z\arabic*},leftmargin=*,resume]
\item \textbf{Balance angles.} Prove uniqueness, monotonicity, and distribution of \(\theta_{\mathrm{bal}}(\rho)\). \hfill (Problem~\ref{prob:angles})
\item \textbf{Explicit formula.} Derive an explicit identity linking \(E(\theta,s)\) to the primes. \hfill (Problem~\ref{prob:explicit})
\end{enumerate}

\noindent\textbf{Phase 4 -- RH}
\begin{enumerate}[label=\textbf{R\arabic*},leftmargin=*,resume]
\item \textbf{RH via selected series.} Prove that \(-\tfrac12\) is not in the range of \(\operatorname{Re}E(\cdot,s)\) for \(\operatorname{Re}s\neq\tfrac12\). \hfill (Problem~\ref{prob:RHfinal})
\end{enumerate}

\section{Conclusion}
The GDST selected/rejected series programme reframes the Riemann Hypothesis as a concrete property of an explicitly defined Dirichlet series attached to a deterministic greedy algorithm.
While the original transfer‑operator route is closed, the cancellation and balance mechanism opens a new avenue.
The programme is problem‑oriented; it does not claim a proof of the Riemann Hypothesis, but it provides a coherent set of questions whose resolution would significantly advance our understanding of the GDST and its relationship to \(\zeta(s)\).

\begin{thebibliography}{9}
\bibitem{GDSTfound}
V.~Geere et al., \emph{The Greedy Dirichlet Sub‑Sum Transform: Foundations, Corrections, and Open Problems}, 2026.
\bibitem{Geere2026a}
V.~Geere, \emph{On the Correlation Structure of a Greedy Harmonic Decomposition of the Sine Function}, March 2026.
\bibitem{Geere2026b}
V.~Geere, \emph{The Greedy Harmonic Decomposition and the Dirichlet Eta Function}, March 2026.
\bibitem{HSTprogramme}
V.~Geere et al., \emph{A Harmonic Sine Transform and the Riemann Zeta Function}, 2026.
\end{thebibliography}

\end{document}